\begin{table}[H]
\centering
\caption{Principal notation used throughout the paper.}
\label{tab:notation}
\small
\setlength{\tabcolsep}{4pt}
\begin{tabularx}{\textwidth}{lY}
\toprule
Symbol & Meaning \\
\midrule
$n$ & number of features \\
$K$ & number of replicates \\
$X_{i,j}$ & raw score of feature $i$ on replicate $j$ \\
$\bU_i=(U_{i,1},\ldots,U_{i,K})$ & uniformized vector with $U_{i,j}=F_j(X_{i,j})$ \\
$Z_i\in\{0,1\}$ & latent reproducibility class indicator \\
$\pi=\Pr(Z_i=1)$ & reproducible-component proportion \\
$c_0,c_1$ & irreproducible and reproducible-component copula densities \\
$F_j$ & marginal CDF of replicate $j$ \\
$\sigma\in[0,1)$ & Pitman--Yor discount parameter \\
$\alpha>0$ & Pitman--Yor concentration parameter \\
$T,T_n$ & variational truncation level or theoretical sieve truncation level \\
$T_{-i}$ & number of occupied reproducible clusters after removing feature $i$ \\
$n_{m,-i}$ & count in occupied cluster $m$ after removing feature $i$ \\
$\theta_m,t_m$ & continuous parameters and atomic-type indicator of atom $m$ \\
$w_m,V_m$ & PY mixture weight and stick-breaking ratio \\
$z_i\in\mathbb N$ & cluster assignment within the reproducible component \\
$M_j$ & Bernstein basis order for marginal $j$ \\
$S_{i,j}$ & Bernstein basis-membership indicator for score $X_{i,j}$ \\
$H$ & number of auxiliary atoms in Algorithm~\ref{alg:mcmc} \\
$h(\cdot,\cdot)$ & Hellinger distance \\
$\epsilon_n,\delta_n$ & posterior contraction rate and local-idr error scale \\
$\widehat k_\alpha$ & number of rejections selected by the Bayes-mFDR step-up rule \\
\bottomrule
\end{tabularx}
\end{table}
